\section{The discrete-time optimality conditions}\label{sec:qsp:model}

\subsection{The Lagrangian and the discount convention}\label{sec:qsp:model:lagrangian}

Adjoin the three static constraints of Problem~\ref{prob:qsp} with multipliers $\beta^t\lambda_t$,
$\beta^t\lambda_tp^W_t$ and $\beta^t\lambda_t\zeta_t$, exactly as in~\eqref{eq:sp:lagrangian}, and the
five laws of motion with multipliers dated at the period in which the stock arrives. Writing the
costate on a stock $Z_{t+1}$ as $\beta^{t+1}\lambda_{t+1}$ times its price in period-$\left(t+1\right)$
goods, the Lagrangian is
\begin{align}
\mathcal L
={}&\sum_{t=0}^\infty\beta^t\Big\{u\!\left(C_t\right)-v_t\!\left(P_t\right)
\notag\\
&+\lambda_t\Big[F_t-C_t-G_t-C^N_t-C^D_t-c^W_tW_t\Big]
+\lambda_tp^W_t\Big[W_t-R_t+\Phi^I_tG_t-\delta M^K_t-\mathcal N_t\Big]
\notag\\
&+\lambda_t\zeta_t\Big[R_t-\Omega_t\left(\mathcal D_t+\phi^I_tG_t\right)\Big]
+\beta\lambda_{t+1}q_{t+1}\Big[\left(1-\delta\right)K_t+G_t-K_{t+1}\Big]
\notag\\
&+\beta\lambda_{t+1}p^S_{t+1}\Big[S_t+D_t-N_t-S_{t+1}\Big]
+\beta\lambda_{t+1}p^X_{t+1}\Big[X_t+D_t-X_{t+1}\Big]
\notag\\
&-\beta\lambda_{t+1}p^M_{t+1}\Big[\left(1-\delta\right)M^K_t+\Phi^I_tG_t-M^K_{t+1}\Big]
-\beta\lambda_{t+1}p^P_{t+1}\Big[\left(1-\theta_t\right)P_t+\Xi_t-P_{t+1}\Big]\Big\}.
\label{eq:qsp:lagrangian}
\end{align}

\smalltitle{One discount convention does all the work.} Define the gross social return between $t$ and
$t+1$ and the corresponding present-value operator
\begin{align}
1+r_{t+1} &\equiv \frac{\lambda_t}{\beta\lambda_{t+1}},
&
\hat p_t &\equiv \frac{p_{t+1}}{1+r_{t+1}}
\quad\text{for}\quad p\in\left\{q,p^S,p^X,p^P,p^M\right\}.
\label{eq:qsp:discount}
\end{align}
So $\hat p_t$ is the value, in period-$t$ goods, of one unit of the stock delivered at $t+1$. With this
single piece of notation the entire discrete-time system is the continuous-time system of
Section~\ref{sec:sp:opt:foc} with each \emph{stock} price replaced by its present value $\hat p_t$, and
with each costate differential equation replaced by a one-period recursion. The two prices that carry
no stock of their own, $p^W_t$ and $\zeta_t$, and the composite $\Psi_t$ built from them, stay
contemporaneous. That correspondence is worth stating before the conditions rather than after, because
it is the only thing one needs to hold in mind while reading them.

Accordingly, redefine the three shorthands of~\eqref{eq:sp:shorthand} and the virgin displacement value
of~\eqref{eq:sp:sum:V} at their discrete-time values,
\begin{align}
\Psi_t &\equiv F_{R,t}\left(1-\zeta_t\Omega^Y_t\right)+\zeta_t,
&
\sigma_t &\equiv \frac{\Phi^I_tG_t}{R_t},
&
V_t &\equiv \Psi_t-p^W_t+d^W\hat p^P_t .
\label{eq:qsp:shorthand}
\end{align}

\subsection{Conditions for the controls}\label{sec:qsp:model:controls}

Differentiating~\eqref{eq:qsp:lagrangian} and dividing through by $\beta^t\lambda_t$ gives the
following, each stated beside the continuous-time condition it discretises.

\emph{Material intensity}, from~\eqref{eq:sp:sum:Omega}:
\begin{align}
\zeta_t &= \sigma_t\left(p^W_t-\hat p^M_t\right).
\label{eq:qsp:foc:Omega}
\end{align}
The discounting is the only change, and it has a reading: storing a tonne in the capital stock saves
today's waste charge $p^W_t$ in full but incurs the embodied-material liability only from $t+1$, so the
liability enters at its present value.

\emph{Consumption}, from~\eqref{eq:sp:sum:C}:
\begin{align}
\frac{u'\!\left(C_t\right)}{\lambda_t} &= 1+\zeta_t\Omega^C_t ,
\qquad\text{with CRRA}\qquad
\lambda_t = \frac{C_t^{-\eta}}{1+\zeta_t\Omega^C_t}.
\label{eq:qsp:foc:C}
\end{align}

\emph{Gross investment}, from~\eqref{eq:sp:sum:I}:
\begin{align}
\hat q_t &= 1-\Phi^I_t\left(1-\sigma_t\right)\left(p^W_t-\hat p^M_t\right).
\label{eq:qsp:foc:G}
\end{align}
It is $\hat q_t$, the period-$t$ price of a unit of capital delivered at $t+1$, that departs from unity
by the net storage value --- which is the right object, since that is what a unit of gross investment
buys.

\emph{Extraction}, from~\eqref{eq:sp:sum:N}, with the corner branch of~\eqref{eq:sp:sum:N-corner}:
\begin{align}
\Psi_t &= C^N_{N\!,t}\left(1+\zeta_t\Omega^N_t\right)+\hat p^S_t+p^W_t\left(1+\Omega^{N,S}_t\right).
\label{eq:qsp:foc:N}
\end{align}

\emph{Exploration}, from~\eqref{eq:sp:sum:D}, with the corner branch of~\eqref{eq:sp:sum:D-corner}:
\begin{align}
\hat p^S_t+\hat p^X_t &= C^D_{D,t}\left(1+\zeta_t\Omega^D_t\right)+p^W_t\Omega^{D,S}_t .
\label{eq:qsp:foc:D}
\end{align}

\emph{Waste}, from~\eqref{eq:sp:sum:W}, the condition that determines $p^W_t$:
\begin{align}
p^W_t\left(1-\alpha_t\varpi_t\right) &= c^W_t\left(1+\zeta_t\Omega^W_t\right)
+\hat p^P_t\left[1-\varpi_t\left(1-d^W\right)-d^W\alpha_t\varpi_t\right]-\alpha_t\varpi_t\Psi_t .
\label{eq:qsp:foc:W}
\end{align}

\emph{Treatment share}, from~\eqref{eq:sp:sum:varpi}, with
$\mathcal G^\varpi_t\equiv\alpha_t\left(\Psi_t-p^W_t\right)+\hat p^P_t\left[\left(1-d^W\right)+d^W\alpha_t\right]$:
\begin{subequations}
\label{eq:qsp:foc:varpi}
\begin{align}
\varpi_t &= 0
&&\text{if}\quad \mathcal G^\varpi_t\le0,
\\
\mathcal G^\varpi_t &= \frac{\bar c^T}{\left(1+g^T\right)^t}\varpi_t^{\xi^T}\left(1+\zeta_t\Omega^W_t\right)
&&\text{if}\quad 0<\mathcal G^\varpi_t<\frac{\bar c^T}{\left(1+g^T\right)^t}\left(1+\zeta_t\Omega^W_t\right),
\\
\varpi_t &= 1
&&\text{if}\quad \mathcal G^\varpi_t\ge\frac{\bar c^T}{\left(1+g^T\right)^t}\left(1+\zeta_t\Omega^W_t\right).
\end{align}
\end{subequations}

\emph{Recycling capital}, from~\eqref{eq:sp:sum:KR}, using $a'_t\!\left(0^+\right)=a^m_t$
and $\mathcal G^K_t\equiv\Psi_t-p^W_t+d^W\hat p^P_t=V_t$:
\begin{subequations}
\label{eq:qsp:foc:KR}
\begin{align}
K^R_t &= 0
&&\text{if}\quad a^m_tV_t\le F_{K,t}\left(1-\zeta_t\Omega^Y_t\right),
\\
\frac{a^m_t}{\left(1+x_t\right)^2}\,V_t &= F_{K,t}\left(1-\zeta_t\Omega^Y_t\right)
&&\text{if}\quad a^m_tV_t> F_{K,t}\left(1-\zeta_t\Omega^Y_t\right),
\end{align}
\end{subequations}
the interior branch of which can be solved for $x_t$ in closed form under this parameterisation,
\begin{align}
x_t &= \sqrt{\frac{a^m_tV_t}{F_{K,t}\left(1-\zeta_t\Omega^Y_t\right)}}-1 ,
\label{eq:qsp:x-closed-form}
\end{align}
which is a genuine convenience: the capital margin never has to be solved numerically, and the corner
test is just the requirement that the square root exceed one.

\subsection{Costate recursions}\label{sec:qsp:model:costates}

Each differential equation of Section~\ref{sec:sp:opt:foc} becomes a one-period recursion of the form
\emph{price today $=$ dividend today $+$ discounted price tomorrow}, with the appropriate survival
factor.

\emph{Capital}, from~\eqref{eq:sp:sum:K}:
\begin{align}
q_t &= F_{K,t}\left(1-\zeta_t\Omega^Y_t\right)+\left(1-\delta\right)\hat q_t,
&&\text{hence}&
1+r_t &= \frac{F_{K,t}\left(1-\zeta_t\Omega^Y_t\right)+\left(1-\delta\right)\hat q_t}{\hat q_{t-1}} .
\label{eq:qsp:costate:K}
\end{align}
The second form is the arbitrage condition: a unit of capital bought at $t-1$ for $\hat q_{t-1}$ pays a
marginal product net of embodied material at $t$ and can be resold, depreciated, for $\left(1-\delta\right)\hat q_t$.

\emph{Reserves and discoveries}, from~\eqref{eq:sp:sum:S} and~\eqref{eq:sp:sum:X}:
\begin{align}
p^S_t &= \left|C^N_{S,t}\right|\left(1+\zeta_t\Omega^N_t\right)+\hat p^S_t,
&
p^X_t &= -C^D_{X,t}\left(1+\zeta_t\Omega^D_t\right)+\hat p^X_t\;<\;0 .
\label{eq:qsp:costate:SX}
\end{align}

\emph{Embodied materials}, from~\eqref{eq:sp:sum:M}:
\begin{align}
p^M_t &= \delta p^W_t+\left(1-\delta\right)\hat p^M_t
\qquad\Longrightarrow\qquad
p^M_t = \sum_{s\ge t}\delta\left(1-\delta\right)^{s-t}p^W_s\prod_{i=t+1}^{s}\frac{1}{1+r_i},
\label{eq:qsp:costate:M}
\end{align}
so $p^M_t$ remains a forward-weighted average of the waste price, with the surviving fraction
$\left(1-\delta\right)^{s-t}$ of the vintage doing the weighting.

\emph{Pollution}, from~\eqref{eq:sp:sum:P}, with
$\mathcal V_t\equiv v'_t\!\left(P_t\right)/\lambda_t-F_{P,t}\left(1-\zeta_t\Omega^Y_t\right)>0$ and
$\hat\theta_t=\left(1-\xi^P\right)\theta_t$ from~\eqref{eq:qsp:absorption}:
\begin{align}
p^P_t &= \mathcal V_t+\left(1-\hat\theta_t\right)\hat p^P_t,
&
\mathcal V_t &= \frac{\psi_tP_t^{\nu}}{\lambda_t}+\gamma\frac{Y_t}{P_t}\left(1-\zeta_t\Omega^Y_t\right),
\label{eq:qsp:costate:P}
\end{align}
the second expression substituting the functional forms of Section~\ref{sec:qsp:setup:forms}. Since
$0<\hat\theta_t<1$ under~\eqref{eq:qsp:absorption} and $r_t>0$, the recursion is a contraction and
$p^P_t>0$ converges, which is the discrete-time counterpart of the remark that this parameterisation
rules out the tipping point.

\emph{Consumption}, from~\eqref{eq:sp:sum:Cgrowth}, combining~\eqref{eq:qsp:foc:C}
with~\eqref{eq:qsp:discount}:
\begin{align}
\frac{C_t^{-\eta}}{1+\zeta_t\Omega^C_t} &= \beta\left(1+r_{t+1}\right)\frac{C_{t+1}^{-\eta}}{1+\zeta_{t+1}\Omega^C_{t+1}} .
\label{eq:qsp:euler}
\end{align}
This is the Euler equation the solver will actually iterate on. A material wedge that is constant
across periods cancels; only a \emph{changing} wedge tilts the consumption path, as in continuous time.

\subsection{The within-period block}\label{sec:qsp:model:block}

Discretisation introduces one object with no continuous-time counterpart, and it is worth isolating
because it is where an implementation is most likely to go wrong. In continuous time the accounting
identities relate flows at an instant and there is nothing to solve. In discrete time, recycled
material re-enters production in the same period in which the waste was generated, so
\eqref{eq:qsp:goods}--\eqref{eq:qsp:intensity} form a simultaneous system inside the period.

Take the states $\left(K_t,S_t,X_t,P_t,M^K_t\right)$ and the five substantive controls
$\left(C_t,N_t,D_t,\varpi_t,K^R_t\right)$ as given. Then $C^N_t$, $C^D_t$, $c^W_t$, $\mathcal N_t$ and
$\phi^I_t$ are known numbers, and the remaining objects are pinned by
\begin{align}
R_t &= N_t+a^m_t\,\frac{\varpi_tW_tK^R_t}{\varpi_tW_t+K^R_t},
&
Y_t &= A_t\left(K_t-K^R_t\right)^{\chi}\left(R_t-\bar R\right)^{1-\chi},
\label{eq:qsp:block-defs}
\\
\mathcal D_t &= \omega^Y_tY_t+\omega^N_tC^N_t+\omega^D_tC^D_t+\omega^W_tc^W_tW_t+\omega^C_tC_t,
&
\sigma_t &= \frac{\phi^I_tG_t}{\mathcal D_t+\phi^I_tG_t},
\label{eq:qsp:block-D}
\end{align}
together with the two equations
\begin{align}
G_t &= Y_t-C_t-C^N_t-C^D_t-c^W_tW_t,
&
W_t &= \left(1-\sigma_t\right)R_t+\delta M^K_t+\mathcal N_t .
\label{eq:qsp:block-system}
\end{align}
The second of~\eqref{eq:qsp:block-system} is the ledger with the intensity condition substituted, using
$\Phi^I_tG_t=\sigma_tR_t$ from~\eqref{eq:qsp:intensity}; $\Omega_t=R_t/\left(\mathcal D_t+\phi^I_tG_t\right)$
is then read off residually.

\begin{lemma}[The within-period block is a well-posed scalar problem]\label{lem:qsp:block}
Given the states and the five controls, \eqref{eq:qsp:block-defs}--\eqref{eq:qsp:block-system} reduce to
a single equation in $W_t$, and any solution satisfies $\Omega_t\ge0$ and $\Phi^Y_t\ge0$
automatically.
\end{lemma}

\begin{proof}
Substituting the first of~\eqref{eq:qsp:block-system} for $G_t$ into the second leaves one equation in
$W_t$, since $R_t$, $Y_t$, $\mathcal D_t$ and $\sigma_t$ are all functions of $W_t$ alone once $G_t$ is
eliminated. The non-negativity claims are~\eqref{eq:sp:setup:Omega-phi}: with $\Phi^I=\Omega\phi^I$ both
hold by construction.
\end{proof}

We stop short of claiming existence and uniqueness, because they are not free. The map has the right
local behaviour: $\partial R_t/\partial W_t=\alpha_t\varpi_t<1$, since
$\alpha_t<a_t\le a^m_t<1$ and $\varpi_t\le1$, so the material feedback from waste back into production
is a contraction --- this is the circularity multiplier $\left(1-\alpha\varpi\right)^{-1}$ seen as a
convergence rate. But $Y_t$ also rises with $W_t$ through $R_t$, and $\mathcal D_t$ with it, so
monotonicity of the whole system is not guaranteed without further restrictions. In practice a scalar
bracketing method on $W_t$ is the natural approach, initialised at
$W_t=N_t+\delta M^K_t+\mathcal N_t$, the value the ledger returns with neither recycling nor material
storage, which bounds the solution from above in the relevant range. \emph{The behaviour of this block,
and whether it stays well conditioned when $\bar R>0$ and $R_t$ approaches the floor, is the first
thing the numerical work should establish.}

\subsection{Solving the model}\label{sec:qsp:model:solution}

Three features of the model determine what solution methods are available, and they rule out the
obvious first choice.

\smalltitle{There is no stationary balanced growth path.} Output and consumption grow with $g^A$, but
the material variables do not: $N_t$ is bounded by the reserve, $W_t$ and $R_t$ are tonnages, and the
whole point of the model is that the reserve is exhausted while the recycling frontier advances. The
economy therefore has no steady state to linearise around and no growth transformation that makes every
variable stationary at once. The usual toolkit of perturbation around a balanced path is unavailable.

\smalltitle{The path has kinks, not smooth transitions.} Proposition~\ref{prop:sp:phases} says the
optimal path switches from $K^R_t=0$ to $K^R_t>0$ at an endogenous date $t_R$, with $\dot K^R$ and
$\dot\varpi$ discontinuous there, and~\eqref{eq:qsp:foc:varpi} adds a second possible switch at
$\varpi_t=1$. Any method that assumes differentiability of the policy function across the whole state
space will handle the transition badly, and the transition is the object of interest.

\smalltitle{The problem is deterministic with a known initial condition.} This is the redeeming
feature. There is no expectation operator anywhere in Problem~\ref{prob:qsp}, so the solution is a
single trajectory rather than a policy function, and the natural formulation is a finite-horizon
boundary-value problem: choose $\left\{C_t,N_t,D_t,\varpi_t,K^R_t\right\}_{t=0}^{T}$ to maximise the
truncated objective subject to the constraints and to terminal conditions at $T$, with $T$ chosen large
enough that the terminal treatment does not contaminate the transition. Two implementations suggest
themselves and should be compared:
\begin{enumerate}[label=\emph{(\alph*)},leftmargin=2.2em,itemsep=0.15em]
\item \emph{Direct nonlinear programming.} Treat the whole trajectory as one large constrained
optimisation, with the within-period block of Section~\ref{sec:qsp:model:block} imposed as equality
constraints rather than solved inside a loop, and the bounds $\varpi_t\in\left[0,1\right]$ and
$K^R_t\ge0$ handled by the solver's own complementarity machinery. This has the advantage of never
requiring the corner branches to be tested by hand, which given two bounded controls and a third corner
at $D_t=0$ is a real advantage.
\item \emph{Shooting on the costates.} Guess the four free initial costates, integrate the FOCs and the
laws of motion forward, and adjust the guess to hit the terminal conditions. Cheaper per iteration and
gives the shadow prices directly --- which matters here, since Part~\ref{part:qmkt} needs
$\left(p^P_t,p^W_t,\zeta_t\right)$ as the policy rates --- but the phase switches make the shooting map
non-smooth.
\end{enumerate}
Method~(a) is the safer starting point and method~(b) is the one that will scale to the policy
experiments, so the sensible plan is to build~(a) first and use it to validate~(b). The comparison is
the first task of the numerical phase.

\smalltitle{Terminal conditions.} The transversality conditions of Section~\ref{sec:sp:opt:foc} have to
be replaced by something implementable at $T$. The natural choices are $q_T$, $p^S_T$, $p^X_T$, $p^M_T$
and $p^P_T$ set to the values implied by treating period $T$ as permanent, that is by solving the
recursions of Section~\ref{sec:qsp:model:costates} at their fixed points with $r_{T+1}=r_T$; and the
horizon should be extended until the transition date $t_R$ and the welfare measure of
Section~\ref{sec:qmkt:welfare} are insensitive to $T$. Given that the reserve is exhausted in finite
time under exploration ceasing --- which~\eqref{eq:sp:sum:D-corner} guarantees --- a horizon of several
centuries at annual frequency is likely to be needed, and the period length is itself a modelling
choice worth testing.
